% lecture01.tex — Conditional Expectation

\section{Topics}
\begin{enumerate}
    \item Conditional Expectation.
    \item Radon--Nikodym Theorem.
\end{enumerate}

\section{Course Overview and Goals}

\textbf{Q.} What is this course about?

\textit{Stochastic Processes} $\{X_t : t \in \T\}$, $X_t \in \cS$ --- state space, $\T$ --- index set.

\textbf{Q.} Why care about stochastic processes? --- \textit{To model \& predict dynamics occurring in the physical world, e.g.\ weather dynamics, infection counts, stock prices, etc.}

\textit{Interest: Can we predict future behavior? Can we use dynamics to produce ``samples'' of interest? --- applications in Gen AI.}

\textit{Our approach in this course:} Focus on Brownian Motion \& allied stochastic processes.

\textit{Why? BM is a canonical stochastic process with many fascinating properties. Using It\^o calculus, we can construct many other stochastic properties from BM.}

\textit{Finally: The techniques \& principles we learn can be widely used to study other stochastic processes.}

\textit{Before diving into a systematic study of BM, we have to build up some technical foundations.}

\section{Review (and Generalization) of Conditional Expectation}

\subsection{Simple Definition}

\begin{definition}{Conditional Expectation (Simple)}{cond-exp-simple}
    Let $X, Y$ be random variables with $\E\abs{Y} < \infty$. The \emph{conditional expectation} $\E[Y \mid X]$ is the (almost surely) unique function $g(X)$ that uncorrelates $Y - g(X)$ from $h(X)$ for all bounded, measurable $h \colon \R \to \R$. Formally,
    \[
        \E\bigl[(Y - g(X)) h(X)\bigr] = 0.
    \]
\end{definition}

\subsection{Limitations}

\begin{enumerate}[label=(\roman*)]
    \item The previous definition works if $X \in \R^n$ for some $n \geq 1$. In working with continuous-time stochastic processes, we want to define quantities like $\E\bigl[X_{t+h} \mid \{X_s : s \leq t\}\bigr]$.

    \textbf{Q.} How to make sense of this?

    \item Existence \& uniqueness of stochastic conditional expectations? (Can use $L^2$-projection if $\E Y^2 < \infty$, but general construction is more involved.)
\end{enumerate}

\subsection{General Definition of Conditional Expectation}

\begin{definition}{Conditional Expectation (General)}{cond-exp-general}
    Let $(\Omega, \F, \Prob)$ be a probability space. $X \colon (\Omega, \F) \to (\R, \B_\R)$ is measurable and $\E\abs{X} < \infty$. Let $\G \subseteq \F$ be a sub-$\sigma$-field.

    Then $\E[X \mid \G]$ is the unique (up to measure zero) $\G$-measurable function such that
    \[
        \E[X \ind_A] = \E\bigl[\E[X \mid \G] \ind_A\bigr] \quad \forall\, A \in \G.
    \]
\end{definition}

\begin{remark}{}{general-def-remarks}
    \begin{enumerate}[label=(\arabic*)]
        \item This definition generalizes the simple definition (Definition~\ref{def:cond-exp-simple}), once we take $\G = \sigma(X)$.
        \item For continuous-time stochastic processes, we can work with
        \[
            \E\bigl[X_{t+h} \mid \sigma(\{X_s : s \leq t\})\bigr].
        \]
    \end{enumerate}
\end{remark}

\textbf{Q.} Existence \& uniqueness? --- \textit{We will come back to this.}

\section{Properties of Conditional Expectation}

\begin{theorem}{Properties of Conditional Expectation}{cond-exp-properties}
    Let $X, Y$ satisfy $\E\abs{X}, \E\abs{Y} < \infty$.
    \begin{enumerate}[label=(\alph*)]
        \item \textbf{(Linearity)} $\E[aX + Y \mid \G] = a\,\E[X \mid \G] + \E[Y \mid \G]$.
        \item \textbf{(Monotonicity)} $X \leq Y \implies \E[X \mid \G] \leq \E[Y \mid \G]$.
        \item \textbf{(Monotone Convergence)} $X_n \geq 0$, $X_n \uparrow X$, $\E X < \infty$
        \[
            \implies \E[X_n \mid \G] \uparrow \E[X \mid \G].
        \]
    \end{enumerate}
\end{theorem}

\begin{proof}
    See Theorem~4.1.9 (Durrett).
\end{proof}

\begin{theorem}{Jensen's Inequality for Conditional Expectation}{jensen-cond-exp}
    Let $\varphi \colon \R \to \R$ be convex. Then
    \[
        \E[\varphi(X) \mid \G] \geq \varphi\bigl(\E[X \mid \G]\bigr).
    \]
\end{theorem}

\begin{proof}[Proof sketch]
    (See Durrett Theorem~4.1.10 for full proof.)

    Assume $\varphi$ is differentiable and $\E\bigl[\abs{\varphi'(\E[X \mid \G])}\bigr] < \infty$ (or $\abs{\varphi'(x)} \leq 1$ for all $x$). By convexity,
    \[
        \varphi(X) \geq \varphi\bigl(\E[X \mid \G]\bigr) + \varphi'\bigl(\E[X \mid \G]\bigr)\bigl(X - \E[X \mid \G]\bigr).
    \]
    Taking conditional expectations,
    \[
        \E[\varphi(X) \mid \G] \geq \varphi\bigl(\E[X \mid \G]\bigr).
    \]

    \textit{Extending to the general case requires some properties of convex functions.}
\end{proof}

\begin{corollary}{}{jensen-corollaries}
    \begin{enumerate}[label=(\arabic*)]
        \item $\abs{\E[X \mid \G]} \leq \E[\abs{X} \mid \G]$. \quad --- Use Jensen's inequality on $\varphi(x) = \abs{x}$.
        \item $X_1, \ldots, X_n \Lpto{1} X \implies \E[X_n \mid \G] \Lpto{1} \E[X \mid \G]$.
    \end{enumerate}
\end{corollary}

\begin{proof}
    For (2): $X_n \Lpto{1} X \implies \E\bigl[\abs{X - X_n} \mid \G\bigr] \Lpto{1} 0 \implies \E[X_n \mid \G] \Lpto{1} \E[X \mid \G]$.
\end{proof}
